%% Согласно ГОСТ Р 7.0.11-2011:
%% 5.3.3 В заключении диссертации излагают итоги выполненного исследования,
%% рекомендации, перспективы дальнейшей разработки темы.
%% 9.2.3 В заключении автореферата диссертации излагают итоги данного
%% исследования, рекомендации и перспективы дальнейшей разработки темы.

\chapter*{Заключение}
\addcontentsline{toc}{chapter}{Заключение}

\begin{vTwoBlock}
В диссертационной работе разработана и апробирована численно-аналитическая схема построения начального приближения для прикладных задач баллистического проектирования межпланетных перелётов с малой тягой. Схема основана на использовании модельного многообразия решений задачи встречи с идеально регулируемой тягой, конечных формул для его аппроксимации и линейных преобразований начальных значений сопряжённых переменных.

Основные результаты работы образуют следующую последовательность.
\begin{enumerate}
	\item Для заданной даты старта по цепочке конечных формул выбирается точка оптимальной дальности в расслоении решений задачи встречи с идеально регулируемой тягой. Этот выбор задаёт исходную модельную траекторию в выбранном семействе решений.

	\item В точке оптимальной дальности по конечным формулам вычисляются продолжительность компланарного перелёта между круговыми орбитами и начальные значения сопряжённых переменных. Тем самым формируется начальное приближение в плоской круговой модельной постановке.

	\item По конечным формулам вычисляются управляющие параметры композиции линейных преобразований начальных значений сопряжённых переменных с учётом коррекции витковой дальности. Эти преобразования обеспечивают переход от модельной постановки к задаче с изменённым начальным фазовым вектором, эллиптичностью и пространственной ориентацией конечной орбиты.

	\item Полученное начальное приближение численно уточняется до решения целевой задачи встречи с идеально регулируемой тягой. При этом начальное приближение сохраняет связь с выбранным семейством решений на модельном многообразии.

	\item Полученное решение задачи с идеально регулируемой тягой продолжается к постановке с дополнительным импульсом скорости в начальный момент времени с учётом ограничений разгонного блока. Это позволяет использовать построенное решение в более прикладной баллистической постановке.

	\item Решение с дополнительным импульсом скорости в начальный момент времени продолжается к постановке с постоянной скоростью истечения. Тем самым предложенная схема связывает модельную задачу с идеально регулируемой тягой с более реалистичной постановкой движения космического аппарата с электрореактивной двигательной установкой.
\end{enumerate}

Основной результат диссертации --- вычислительная схема, позволяющая через модельное многообразие, конечные формулы и линейные преобразования начальных значений сопряжённых переменных перейти от расслоения решений задачи встречи с идеально регулируемой тягой к начальному приближению для прикладной баллистической задачи.

Дальнейшее развитие работы связано с переносом предложенной вычислительной схемы на более сложные постановки задачи оптимизации межпланетных перелётов, выходящие за рамки модели идеально регулируемой тяги, и на задачи с более общими граничными условиями.
\end{vTwoBlock}